\honest{Brain Breakthrough! It's Made of Neurons!}{Eliezer Yudkowsky}{2008}

Amazing breakthrough: a team led by the Nobel laureate Santiago Ram\'on y Cajal announces that the brain is a ``ridiculously'' complicated network of tiny cells. The team includes Leeuwenhoek and ``possibly Imhotep''. \nb{This is the Amazing Breakthrough Day the author proposed a week earlier: report old science as news on April 1st.}

Their statement: years of research show that the squishy thing in our skulls is even more complicated than it looks; thanks to Cajal's use of Golgi's new stain, we know the brain is not a continuous network but is made of many tiny cells. \nb{Golgi used the stain to argue that the brain is a continuous network, and held to that view in the Nobel lecture.} From Alcmaeon to Broca, the evidence says the brain is the seat of reason. Nemesius argued that brain tissue was too earthy to link body and soul, and put the mental faculties in the ventricles; but then there is no reason for the brain to be so intricate. Babbage suggested that many small mechanical devices could form an Analytical Engine able to do arithmetic, which is widely thought to require thought. Galvani and Helmholtz show that neurons work electrochemically, not by mechanical pressure, but the analogy still suggests that a vast network of neurons could think.

``We have found an enormously complicated material system located where the mind should be.'' So Spinoza was right and Descartes wrong: ``Mind and body are of one substance.'' \nb{For Spinoza, on the reading the \textsc{Stanford Encyclopedia of Philosophy} adopts, thought and extension have nothing in common and do not act on each other, and thought is as basic as matter. The release's next conclusion, that intelligence is not fundamental, is not Spinoza's.} Add Darwin, and the evidence ``seems to indicate that intelligence is ontologically non-fundamental and has an extended origin in time,'' which ``strongly weighs against'' every theory that makes the mind fundamental, ``including all religions ever invented.'' \nb{The religions are not discussed. And the evidence did not settle the matter for everyone who knew it best: John Eccles, who shared the 1963 Nobel Prize for work on the synapse, remained a dualist.}

Much remains to be done, but we offer ``the promise, though not yet the realization, of a full scientific account of thought.'' The problem may now be declared, if not solved, then solvable.

We regret that Cajal and most of his colleagues are no longer available for comment.
